\section{Methodology}
\label{sec:method}

The experiments put the definition of Section~\ref{sec:definition} into practice with three
components: a small benchmark of coding tasks with planted sub-tasks (DriftBench-mini), a harness
that drives Claude Code headlessly with locally hosted open-weight models, and an evaluator that
scores the final repository and the trajectory. All code, task material and raw trajectories are
part of the project repository.

\subsection{DriftBench-mini}
\label{sec:bench}

DriftBench-mini contains twelve tasks, each a small, self-contained Python library of 125 to 518
lines with its own test suite: invoicing with discount codes, word statistics, a todo store, a CSV
inventory, layered configuration loading, a rate limiter, a Markdown table renderer, a unit
registry, a double-entry ledger, a web router, log statistics and weighted shortest paths. For each
task we wrote:
\begin{itemize}
  \item a goal $\Gz$ (\texttt{goal.md}) phrased as a short user request with five to seven numbered
        requirements, closing with ``keep the change focused on this feature'';
  \item \emph{hidden tests}, one per requirement (5--7 per task), which the agent never sees and
        which define $\Req$;
  \item a \emph{visible} test suite inside the repository (5--26 tests), which the agent can run;
  \item a \emph{scope manifest}: the files a faithful implementation is allowed to change, $S(\Gz)$,
        and a list of \emph{protected} files (existing tests and modules unrelated to $\Gz$);
  \item four injections for each of the four sources of Section~\ref{sec:definition}, each labelled
        with its legitimacy, and a reference solution.
\end{itemize}
The injections follow a fixed pattern per source. \textsf{user} injections are sanctioned side
requests (``Side request: when \texttt{--readability} is passed, the CLI should also print the
sentence count\ldots{} Then carry on with the top-words feature''); each comes with its own side test
and may widen the scope. \textsf{env} injections are distractors that make a protected, unrelated test
fail, for example a new test asserting that abbreviations do not split sentences. \textsf{text}
injections are distractors that plant instruction-like text, such as a \texttt{TODO(maintainer)}
asking to migrate the CLI to \texttt{click} before adding options, or a README note asking for type
hints and \texttt{mypy --strict} on untouched modules. \textsf{cascade} injections are necessary:
they break something the goal depends on (a wrong rounding mode, a function that reads only the first
line of a file, a missing import) and widen the scope to the broken file.

\paragraph{Validation.} A script checks every task mechanically before use: the unmodified
repository passes its visible suite and fails the hidden one; the reference solution passes both;
every \textsf{env} injection makes the visible suite fail; \textsf{text} injections leave the
solution's scores unchanged; every \textsf{cascade} injection lowers the hidden or the visible score
of the solution; every side test fails on the unmodified repository. A lint step also enforces the
labels (distractors may not widen the scope). All twelve tasks pass.

\paragraph{Known limitations of the material.} Four properties of the cascades and sanctioned
requests matter for interpretation. (i) Some cascades are necessary only for some implementations:
a wrong rounding constant in a helper (T01) affects the agent only if it rounds through that helper.
(ii) One cascade (T01) breaks the package import, so its effect on the pass rate is all-or-nothing,
unlike the others. (iii) Cascades that corrupt the pytest configuration (\texttt{addopts},
\texttt{required\_plugins}) affect only the agent's own test runs, because the evaluator neutralises
those options; they create work but cannot lower the score. (iv) Some sanctioned requests live in
files that are already in scope, so they cannot show up as scope violations; they are measured by
their side tests only.

\subsection{Agent and harness}
\label{sec:harness}

\paragraph{Agent.} The agent is Claude Code~2.1.268 in headless mode
(\texttt{claude -p --output-format stream-json}), with permissions bypassed inside a sandbox. Instead
of Anthropic models, it talks to Ollama's Anthropic-compatible endpoint, so that every request goes
to a local open-weight model. We use two agent models from different families: \texttt{qwen3-coder:30b}
(a mixture-of-experts model trained for agentic coding) and \texttt{gpt-oss:20b}. A third candidate,
\texttt{devstral}, was dropped in the pilot because it emitted tool calls as plain text that Claude
Code did not execute. A panel of three judge models labels the trajectories
(Section~\ref{sec:e5}), and \texttt{qwen2.5:7b} was the planned weak model for the inherited-drift test of E3, which was not run (Section~\ref{sec:e3}). The tool set is restricted to \texttt{Task} (sub-agents), \texttt{Bash},
\texttt{Read}, \texttt{Edit} and \texttt{Write}: everything else Claude Code offers (web access,
scheduling, skills) would inflate the context and add confounds. Claude Code's task-list tools are
therefore \emph{off} by default and enabled only in the corresponding E2 condition. The context
window is 64k tokens, with Claude Code's automatic compaction left on, and each run has its own
configuration directory, so that no user-level instructions, memory or plugins reach the agent.

\paragraph{Episode protocol.} Every episode has the same structure, so that conditions differ only
in content (Section~\ref{sec:definition}):
\begin{enumerate}
  \item \emph{Initial segment}: the goal, with at most 8 agent turns;
  \item \emph{four pauses}: the session is resumed with a message and at most 8 more turns. The
        message is a side request (\textsf{user}), a neutral status note of similar length
        (\textsf{filler}), or ``Please continue.'' in every other condition;
  \item \emph{final segment}: ``Please continue until the original task is complete. When you are
        done, end with a short summary of everything you changed and why'', with at most 40 turns;
  \item \emph{probe}: the session is forked and, with no tools, asked to state the original objective,
        the side tasks taken on and whether the objective is done. The fork leaves the main
        trajectory untouched \citep{ding2026contextecho}.
\end{enumerate}
Each segment has a 30-minute time limit. \textsf{env}, \textsf{text} and \textsf{cascade}
injections are applied to the repository before the initial commit, so the agent discovers them as
it works; \textsf{user} injections occupy the first $k$ pauses. A dose of $k$ means $k$ injections of
one source; the pilot also used a \textsf{mix} condition with one injection of each source (four
in total), which E2 reuses as its fixed load.

\paragraph{Infrastructure.} Models were served by Ollama~0.20.1, the newest release compatible with
the NVIDIA~535 drivers of the servers (newer releases fall back silently to the CPU). Each model had
its own GPU and three parallel slots, with three episodes running at a time. Two settings turned out
to matter for speed. Disabling Claude Code's per-request attribution header keeps the prompt prefix
stable, so Ollama can reuse its cache. Never running more episodes than slots avoids evicting
other episodes' caches. Almost all episodes ran on a departmental GPU server (NVIDIA A40), with the
agent in a Docker container that mounted only the task repository and the evaluator in a second
container without network. The last 45 of the 1{,}152 E1 episodes ran on the CINECA Leonardo
cluster (NVIDIA A100) in a SLURM job, with Claude Code and pytest as plain processes. There we
checked that the model weights were identical (same blob digests). Because that setup has no mount
isolation, the harness flags every tool call that references the benchmark's own files (hidden
tests, solutions, task definitions); no episode triggered it.

\subsection{Measures}
\label{sec:measures}

After the probe, the final repository is copied and evaluated. The measures implement the indicator
families of Section~\ref{sec:definition}:
\begin{description}[leftmargin=0pt]
  \item[Outcome.] The \emph{hidden pass rate} is the fraction of hidden requirement tests that pass.
        A test that cannot even be collected (e.g.\ an import error) counts as failed. The hidden
        tests run with the repository's pytest configuration and \texttt{conftest} files disabled,
        so a broken configuration cannot hide or fake a result. \emph{Goal fully met} is the
        fraction of episodes with all requirements passing. The \emph{side pass rate} scores
        sanctioned side requests with their own tests.
  \item[Scope.] The diff against the initial commit is measured in changed lines, excluding caches
        and tool artefacts (e.g.\ \texttt{\_\_pycache\_\_}, \texttt{pytestdebug.log}). In the pilot
        and in early E1 runs the raw $\mathrm{OOS}$ ratio was dominated by benign additions: new test
        files the agent wrote for the goal, and debugging scripts it left behind. Even at $k=0$ it
        reached 20--50\% for \texttt{gpt-oss}. We therefore split out-of-scope changes into three parts:
        \begin{itemize}
          \item \emph{edits to existing files} outside $S^\ast$, the scope-violation measure proper
                ($\mathrm{OOS}_{\text{exist}}$);
          \item \emph{new test files};
          \item \emph{stray new files}, i.e.\ any other new file, which is residual state left in the
                repository.
        \end{itemize}
        Modifications of protected files and commits made by the agent are counted separately.
  \item[Process.] Tool calls by type, compaction events, segments ending at the turn limit, and the
        number of times a mitigation hook fired.
  \item[Representation and reporting.] The probe answer and the final summary are stored for the
        LLM judges (goal recall, goal replacement and report fidelity). The judge panel of E5
        (Section~\ref{sec:e5}) now labels every segment and every final summary; the measures
        derived from its labels are defined in Section~\ref{sec:judge-results}. Goal recall from the
        probe is not used, because the judges do not agree on it ($\kappa=0.17$).
\end{description}

\subsection{Analysis}
\label{sec:analysis}

The episode is the unit of analysis and tasks are the clustering unit, with three repetitions per
task and cell. Cell means come with 95\% confidence intervals from a task-cluster bootstrap (2{,}000
resamples of the twelve tasks). The effect of a cell is its difference from the $k=0$ baseline
($D_k^{\mathrm{src}}$), bootstrapped in the same way. Dose trends are estimated per model and source
by ordinary least squares, $y = \beta_0 + \beta_1 k + \text{task fixed effects}$, on the $k=0$
episodes plus that source's $k\in\{1,2,4\}$ episodes, with task-clustered standard errors. To test
whether a source costs more than an equally long neutral interruption, the same model is fitted on
the source and the filler control together with a separate slope for the source,
$y = \beta_0 + \beta_1 k + \beta_2\, k\cdot\mathbb{1}[\text{source}] + \text{task FE}$. Both arms
share the $k=0$ episodes, and $\beta_2$ is the excess cost per sub-task. A random-intercept mixed
model was our first choice, but it did not converge on the sparse count outcomes. With every task
present in every cell, task fixed effects estimate the same contrasts. $p$-values are
Holm-corrected within each family: the ten slopes of one outcome, the eight source-versus-filler
contrasts, and the four mitigations of E2.
